\begin{description}
	\item[1行目] 各群のサンプルサイズNを設定します。各群共通のサイズになります。
	\item[2行目] 水準数です。ここでは5群の平均値の比較をすることにしました。
	\item[3行目] 全体平均の設定です。事前分布に[0,100]の一様分布を置いてますから，真ん中ぐらいにしてあげました。
	\item[4行目] 誤差の散らばりを設定します。
	\item[6行目] 水準数$-1$の効果を設定します。手入力でもいいのですが面倒なので，-10から10までの範囲で一様乱数によって生成しました\footnote{一様乱数の関数は，Rでは\texttt{runif}です。}。
	\item[7行目] 水準数Lvの効果にするため，先ほど作った水準数$-1$の\verb|raw_effect|に\verb|0-sum(raw_effect)|の計算結果をつけくわえています。関数\texttt{c}は結合させるcombineという意味です。
	\item[8行目] 各群の平均値です。全体平均に効果を足しています。全体平均は1つの数字，効果は水準数の要素を持つベクトルですから，計算ができないように思えますが，このような場合Rは自動的にサイズ合わせのため値の再利用を行います。つまり，$50+(\delta_1,\delta_2,\delta_3,\delta_4,\delta_5)$を$(50,50,50,50,50)+(\delta_1,\delta_2,\delta_3,\delta_4,\delta_5)$と解釈してくれるのです。
	\item[10行目] 水準数$\times$各群のサイズの乱数を発生させています。乱数は正規分布によるもので，平均やSDは上で指定した通りです。ここでも値の再利用が行われていて，ベクトル\verb|mu|は$\mu_1,\mu_2,\mu_3,\mu_4,\mu_5,\mu_1,\mu_2,....$とくり返して当てはめられていきます。
	\item[11-14行目] 数値を\verb|data.frame|型にくみ上げていってます。\verb|rep|関数は，指定した回数だけ繰り返すもので，ここでは\verb|1:Lv|すなわち1,2,3,4,5という数字をデータサイズ分繰り返していることになります。
	\item[16行目] ここでは\verb|cmdstanr|パッケージを使ってコンパイルしています。
	\item[17行目] 推定に使うデータセットです。データ長は\verb|data.frame|の長さを返す関数\verb|NROW|を使って与えています。インデックス変数は\verb|dat|オブジェクトの\verb|Idx|変数，データは同じオブジェクトの\verb|value|変数を指定しています。
\end{description}
